%% kapa: exam=1 round=25 year=2014 period=2 form=A subject=accounting subject_ko=회계학 seq=13 number=53
%% source: Q-Net 감정평가사 기출문제 (cst00309) artlSeq=5101207 file=extracted/1008/5101207/2014년 제25회 감정평가사 1차시험 기출문제/2교시(A형).hwp
%% conversion: kordoc 4.24.51 -> tools/gen1.py (auto); review=proofread-sonnet; flags=none
\begin{question}{53}
(주)한국은 결산을 앞두고 당좌예금의 계정 잔액을 조정하기 위해 은행에 예금 잔액을 조회한 결과 20×1년 12월 31일 잔액은 ￦125,400이라는 회신을 받았다. (주)한국의 당좌예금 장부상의 수정전 잔액은 ￦149,400이다. (주)한국의 내부감사인은 차이의 원인에 대해 분석하였고, 다음과 같은 사실을 확인하였다.
\begin{framedtext}
○ (주)한국이 20×1년 12월 31일에 입금한 ￦50,000이 은행에서는 20×2년 1월 4일자로 입금 처리되었다. \par
○ (주)한국이 발행한 수표 중에서 20×1년 12월 3일에 발행한 수표(no. 164) ￦20,000이 아직 인출되지 않았다. \par
○ (주)한국이 발행한 수표(no. 173)의 발행액은 ￦21,000이었으나 회계담당자가 이를 ￦12,000으로 잘못 기록하였다. \par
○ (주)한국이 발행한 수표(no. 182) ￦15,000을 은행의 착오로 다른 기업의 계좌에서 출금처리하였다. \par
\end{framedtext}
위 자료를 이용할 때 20×1년 말 (주)한국의 수정후 당좌예금 잔액은?
\begin{choices}
  \item ￦134,400
  \item ￦140,400
  \item ￦158,400
  \item ￦168,400
  \item ￦171,400
\end{choices}
\end{question}
\answer{2}
